% ===========================================================================
%  No Lab Required — typographic engine
%  Implements PART B of workbook-production-spec.md v1.0.
%  Engine: LuaLaTeX (needed for OpenType fonts *and* LaTeX PDF tagging).
% ===========================================================================

% --- 1. PAGE GEOMETRY (spec B1, B2 simplified fallback) --------------------
% One column, 5.05 in wide, sitting slightly above the centre of the sheet.
%
% The spec offers a single column as its own fallback and it is the better
% design here, for a reason that only showed up once the book was written: the
% side column was carrying 62 margin definitions and every one of them
% repeated, nearly word for word, a sentence already in the body two inches to
% its left. A reader met the same definition twice and the page sat lopsided.
%
% The measure is 5.05 in rather than 5.50 because at 11 pt Charis a 5.50 in
% line runs to about 80 characters, and a reader loses the start of the next
% line somewhere past 75. 5.05 in gives roughly 74. Eighty columns of
% Inconsolata at 9 pt is 5.00 in, so code still fits inside the measure and
% nothing has to break out of the column.
%
% Top 1.05 in against bottom 1.45 in puts the block above the optical centre,
% which is where books have put it since before anyone measured why.
\usepackage{geometry}
\geometry{
  paperwidth      = 8.5in,
  paperheight     = 11in,
  top             = 1.05in,
  bottom          = 1.45in,
  left            = 1.725in,
  right           = 1.725in,
  twoside         = true,
  asymmetric      = true,
  headheight      = 12pt,
  headsep         = 38.0pt,  % running-head baseline 0.50in from the top trim
  footskip        = 52.0pt,  % folio baseline 0.72in from the bottom trim
}

\newlength{\nlrfullwidth}   \setlength{\nlrfullwidth}{5.05in}
\newlength{\nlrbreakout}    \setlength{\nlrbreakout}{0pt}

% --- 2. FONTS (spec B3) ----------------------------------------------------
% Two families, not three, and neither is the open-source default everybody
% reaches for first.
%
% Charis is Matthew Carter's Charter, extended by SIL. Carter drew Charter to
% hold up at small sizes on the cheap low-resolution printers of the 1980s.
% This book assumes its reader prints on a school printer in black and white,
% so that is not a nostalgic choice, it is the same problem.
%
% Inconsolata was drawn by Raph Levien specifically for printed code listings,
% which is exactly what a code block in a workbook is. It advances 0.5 em per
% character, half the width of most monospaces, so 80 columns fits the measure.
%
% Headings are set in the text face at a heavier weight rather than in a
% separate sans. Books have done that for a long time and it stops the page
% looking like a component library.
\usepackage{fontspec}
% Ligatures=TeX (the alias for the tex-text input mapping) is applied to the
% text faces one at a time rather than to every font in the document. Set as a
% default it reached Inconsolata too, and the mapping turns a straight quote
% into a typographic one, so every string literal in every listing printed as
% "basal" instead of "basal" and no code in the book could be typed off the
% page. There is no Ligatures=NoTeX to undo it with; the fix is not to set it.

\setmainfont{Charis}[
  UprightFont    = *-Regular,
  ItalicFont     = *-Italic,
  BoldFont       = *-Bold,
  BoldItalicFont = *-BoldItalic,
  Ligatures      = TeX,
  Scale          = 1,
]

% Scale=1 is load-bearing, not decoration. The 80-column guarantee is a width
% in inches -- 80 x 0.5 em x 9 pt = 360 pt = 5.00 in against a 5.05 in measure
% -- and it only holds if 9 pt means 9 pt. Pandoc's template sets
% \defaultfontfeatures{Scale=MatchLowercase}, which matches Inconsolata's
% lowercase to Charis's and silently renders the code at 9.46 pt: 80 columns
% becomes 5.25 in and the longest listings run off the text block. This was
% hidden for a while because the old \defaultfontfeatures{Ligatures=TeX} here
% overwrote pandoc's line wholesale, so removing that one exposed it.
% No Ligatures=TeX either, deliberately. See the note above \setmainfont.
\setmonofont{Inconsolata}[
  UprightFont    = *-Regular,
  BoldFont       = *-Bold,
  ItalicFont     = *-Regular,
  Ligatures      = NoCommon,
  Scale          = 1,
]

\newfontfamily\nlrsans{Charis}[Scale=1, UprightFont=*-Regular, ItalicFont=*-Italic, BoldFont=*-Bold, BoldItalicFont=*-BoldItalic, Ligatures=TeX]
\newfontfamily\nlrsansmed{Charis}[Scale=1, UprightFont=*-Regular, ItalicFont=*-Italic, BoldFont=*-Bold, BoldItalicFont=*-BoldItalic, Ligatures=TeX]
\newfontfamily\nlrsanssemi{Charis}[Scale=1, UprightFont=*-Bold, ItalicFont=*-BoldItalic, BoldFont=*-Bold, BoldItalicFont=*-BoldItalic, Ligatures=TeX]
% The title page only. A word space that is proportionally right at 11 pt
% reads loose at 44, and WordSpace is not settable through \addfontfeatures --
% it has to be declared with the family.
\newfontfamily\nlrdisplay{Charis}[UprightFont=*-Bold, ItalicFont=*-BoldItalic,
  BoldFont=*-Bold, BoldItalicFont=*-BoldItalic, Ligatures=TeX, Scale=1,
  WordSpace=0.90]
\newfontfamily\nlrmono{Inconsolata}[UprightFont=*-Regular, BoldFont=*-Bold, ItalicFont=*-Regular, Ligatures=NoCommon, Scale=1]

\newcommand{\nlrlining}{}

% --- 3. COLOUR (spec B5) ---------------------------------------------------
% Ink, one grey, and three muted accents. Every one has to survive being
% converted to grey on a black-and-white printer, which is the constraint that
% keeps them dark. Nothing carries meaning in hue alone; colour always rides
% alongside a label or a change of style.
%
% The four background tints this file used to define sat within 8/255 of one
% another. On the school printer the book is designed for they all collapsed to
% the same flat grey, so the system spent four values to print one.
\usepackage{xcolor}
\definecolor{nlrink}{HTML}{1C1B19}      % body text, slightly warm black
\definecolor{nlrsecond}{HTML}{5F5B54}   % secondary text
\definecolor{nlraccent}{HTML}{7A2E1D}   % oxblood, external addresses only
\definecolor{nlrteal}{HTML}{7A2E1D}     % alias, older rules still name it
\definecolor{nlramber}{HTML}{8A6212}    % dark ochre
\definecolor{nlrgreen}{HTML}{3B5C3F}    % deep green
\definecolor{nlrrule}{HTML}{B9B4AA}     % rules, dark enough to print

% \color at preamble time does not survive. \begin{document} issues
% \normalcolor, which resets to the driver's black, so every page of this book
% printed at #000000 rather than at the warm black defined above. Setting
% \normalcolor itself is what makes the choice stick.
\AtBeginDocument{\renewcommand{\normalcolor}{\color{nlrink}}\normalcolor}

% --- 4. TYPE SCALE (spec B4) -----------------------------------------------
% Body sits on the 16 pt baseline grid. Everything else is expressed as an
% explicit size/leading pair so the grid stays legible even where it is broken.
\newcommand{\nlrbody}{\fontsize{11}{16}\selectfont}
\newcommand{\nlrcodesize}{\fontsize{9}{12.5}\selectfont}
\newcommand{\nlroutsize}{\fontsize{8.5}{12}\selectfont}
\newcommand{\nlrmarginsize}{\fontsize{8.5}{12}\selectfont}
\newcommand{\nlrcapsize}{\fontsize{9.5}{13}\selectfont}
\newcommand{\nlrtablesize}{\fontsize{9.5}{13}\selectfont}
\newcommand{\nlrcalloutsize}{\fontsize{10}{14.5}\selectfont}
\newcommand{\nlrlabelsize}{\fontsize{8.5}{10}\selectfont}
\renewcommand{\normalsize}{\nlrbody}
\normalsize

\usepackage{microtype}
% Every component label is one italic word or phrase in the text face, closed
% with a full stop, run into the first line of what it labels. Not letterspaced
% capitals, not a coloured word, not a word beside an icon. Capitals shout, and
% a page of them makes a book look like a component library rather than like
% something somebody wrote.
\newcommand{\nlrlabel}[1]{{\itshape #1}}
\newcommand{\nlrlabelmed}[1]{{\itshape\color{nlrsecond}\fontsize{9}{12}\selectfont #1}}

% --- 5. COMPOSITION RULES (spec B8) ----------------------------------------
\usepackage{ragged2e}
% Indented paragraphs with no space between them. The alternative, a zero
% indent plus 8 pt of \parskip, is the exact signature of markdown converted to
% PDF, and it was the single most machine-made thing about the running text.
% Order matters, and getting it wrong is silent. ragged2e's \RaggedRight ends
% by assigning \parindent from \RaggedRightParindent (0pt by default) and
% \rightskip from \RaggedRightRightskip, so setting either one afterwards has
% no effect on the paragraphs it governs. Set before, not after: the indent was
% being zeroed on every paragraph in the book, and with \parskip at 0 as well
% there was nothing left to mark where one paragraph ended and the next began.
% 5em, not the 2.2em this was tuned to. Inline code cannot hyphenate, so a
% sentence ending in \texttt{FileNotFoundError:} gives TeX a choice between a
% short line and an overfull one, and with only 2.2em of rag allowed it took
% the overfull line and ran off the text block. The measured right edge across
% the body still has a median of 6.68 in against a 6.775 in margin, so the rag
% reads as even; the extra tolerance is spent on the handful of lines that
% need it.
\setlength{\RaggedRightParindent}{1.2em}
\setlength{\RaggedRightRightskip}{0pt plus 5em}
\RaggedRight
\setlength{\parindent}{1.2em}
\setlength{\parskip}{0pt plus 0.5pt}
\frenchspacing                        % one space after a full stop, as in print
% \rightskip's stretch is finite under RaggedRight, so a line ending in a long
% unbreakable identifier (\texttt{FileNotFoundError:} inside an indented aside)
% could leave TeX a choice between a gap wider than the rag allows and an
% overfull line, and it took the overfull one. \emergencystretch applies only
% on the pass that would otherwise overflow, so the rest of the rag is untouched.
\emergencystretch=2em
\hyphenpenalty=2000
\doublehyphendemerits=100000
\finalhyphendemerits=100000
\lefthyphenmin=3
\righthyphenmin=3
\widowpenalty=10000
\clubpenalty=10000
\displaywidowpenalty=10000
\brokenpenalty=10000
\raggedbottom

% babel and csquotes are already loaded by the Quarto/pandoc template from the
% document's `lang` setting; loading them again here causes an option clash.

% --- 6. FULL WIDTH ---------------------------------------------------------
% Kept as a no-op so the component definitions and the filter do not need to
% know the layout changed. With one centred column there is nothing wider to
% break out to: 80 columns of Inconsolata at 9.5 pt is 5.28 in and the measure
% is 5.50 in, so code already fits.
\usepackage{changepage}
\usepackage{calc}
\newenvironment{nlrfull}{\par}{\par}

% --- 7. HEADINGS (spec B4) -------------------------------------------------
% Every \titleformat in this file used to be dead.
%
% _quarto.yml turns on PDF tagging, which emits \DocumentMetadata{tagging=on},
% and that replaces the sectioning commands with LaTeX's template-based ones.
% titlesec detects the replacement, prints "Non standard sectioning command
% detected. Using default spacing and no format." into a log nobody reads, and
% then does nothing. So the chapter openers printed stock book.cls: the words
% "Chapter 4" in bold black above a bold black title, with the 60 pt numeral
% and the rule this file described existing only in the comment.
%
% The supported route under tagging is to redeclare the heading instances.
%
% What is set here differs from stock on four axes rather than one. Size alone,
% in a single weight and a single colour, is how a generated page builds
% hierarchy; the label line is also italic where the title is roman, grey where
% the title is ink, and light where the title is bold.
\usepackage{needspace}

\makeatletter
\DeclareInstance{heading}{part}{display}{
  name           = part,
  level          = -1,
  placement      = page,
  after-penalty-vspace = 150pt,   % drops the part title to the optical third
  after-vspace   = 40pt,
  prefix         = \partname,
  number-format  = \thepart,
  heading-decls  = \RaggedRight\parindent0pt\itshape\color{nlrsecond}%
                   \fontsize{11}{15}\selectfont,
  title-decls    = \upshape\bfseries\color{nlrink}\fontsize{32}{36}\selectfont,
  headformat-instance = part,
  mark-cmd       = \partmark{#1},
}
\DeclareInstance{headformat}{part}{display}{
  heading-indent    = 0pt,
  prefix-number-sep = \WordSpaceAmount{1},
  number-title-sep  = 16pt,
}

\DeclareInstance{heading}{chapter}{display}{
  name           = chapter,
  level          = 0,
  placement      = top,
  after-penalty-vspace = 50pt,
  after-vspace   = 32pt,
  prefix         = \chaptername,
  number-format  = \thechapter,
  heading-decls  = \RaggedRight\parindent0pt\itshape\color{nlrsecond}%
                   \fontsize{11}{15}\selectfont,
  title-decls    = \upshape\bfseries\color{nlrink}\fontsize{24}{28}\selectfont,
  headformat-instance = chapter,
  mark-cmd       = \chaptermark{#1},
  contents-extra = \addtocontents{lof}{\protect\addvspace{10\p@}}
                   \addtocontents{lot}{\protect\addvspace{10\p@}},
}
\DeclareInstance{headformat}{chapter}{display}{
  heading-indent    = 0pt,
  prefix-number-sep = \WordSpaceAmount{1},
  number-title-sep  = 14pt,
}

\DeclareInstance{heading}{section}{display}{
  name           = section,
  level          = 1,
  before-vspace  = 22pt plus 4pt minus 2pt,
  after-vspace   = 7pt,
  heading-decls  = \RaggedRight\normalfont\bfseries\color{nlrink}%
                   \fontsize{15}{19}\selectfont,
  headformat-instance = section,
  mark-cmd       = \sectionmark{#1},
}
\DeclareInstance{heading}{subsection}{display}{
  name           = subsection,
  parent-name    = section,
  level          = 2,
  before-vspace  = 15pt plus 3pt minus 1pt,
  after-vspace   = 4pt,
  heading-decls  = \RaggedRight\normalfont\bfseries\color{nlrink}%
                   \fontsize{12}{16}\selectfont,
  headformat-instance = subsection,
  mark-cmd       = \subsectionmark{#1},
}
\DeclareInstance{heading}{subsubsection}{display}{
  name           = subsubsection,
  parent-name    = subsection,
  level          = 3,
  before-vspace  = 12pt plus 2pt,
  after-vspace   = 2pt,
  heading-decls  = \RaggedRight\normalfont\itshape\color{nlrink}%
                   \fontsize{11}{15}\selectfont,
  headformat-instance = subsubsection,
  mark-cmd       = \subsubsectionmark{#1},
}
\makeatother

% Headings never sit alone at the foot of a page.
\pretocmd{\section}{\Needspace*{5\baselineskip}}{}{}
\pretocmd{\subsection}{\Needspace*{4\baselineskip}}{}{}

% --- 8. PAGE FURNITURE (spec B6) -------------------------------------------
\usepackage{fancyhdr}
\usepackage{emptypage}

\newcommand{\nlrrunningfont}{\itshape\fontsize{9}{11}\selectfont\color{nlrsecond}}
\newcommand{\nlrfoliofont}{\nlrsansmed\fontsize{9.5}{11}\selectfont\color{nlrink}\nlrlining}

\fancypagestyle{nlrmain}{%
  \fancyhf{}%
  \renewcommand{\headrulewidth}{0pt}%
  \renewcommand{\footrulewidth}{0pt}%
  % Verso carries the book title, recto the current chapter (spec B6). Both hang
  % on the outer edge of the text block, which is why head and foot are widened
  % to the full 6.50 in below.
  \fancyhead[LE]{\nlrrunningfont\nlrbooktitle}%
  \fancyhead[RO]{\nlrrunningfont\nlrchaptermark}%
  % Outer edge on both sides. \fancyfoot[R] is unconditional, so on a verso
  % the folio sat bottom-right while its own running head hung flush left,
  % with the two on opposite edges of the same page.
  \fancyfoot[LE,RO]{\nlrfoliofont\thepage}%
}
\fancypagestyle{nlrplain}{% chapter openers, part titles: no head, no folio
  \fancyhf{}%
  \renewcommand{\headrulewidth}{0pt}%
  \renewcommand{\footrulewidth}{0pt}%
}
\pagestyle{nlrmain}

% Head and foot span the whole text block, not just the main column, so the
% folio and the recto running head land on the outer edge.
\fancyhfoffset[R]{\nlrbreakout}

\newcommand{\nlrbooktitle}{No Lab Required}
\newcommand{\nlrchaptermark}{}
\renewcommand{\chaptermark}[1]{\renewcommand{\nlrchaptermark}{#1}}
\renewcommand{\sectionmark}[1]{}

\let\nlroldchapter\chapter
% Chapter openers use \thispagestyle{plain}, so `plain` has to become this
% book's plain: no running head, no folio.
%
% \makeatletter must wrap the \AtBeginDocument line itself, not sit inside the
% braces. The hook's body is tokenized when this line is READ, and @ is not a
% letter in a header include, so \ps@plain tokenizes as \ps followed by the
% characters @plain. \let then binds \ps to @ and the rest typesets, which put
% the literal text "plain@@nlrplain" on the half title and cost two pages.
\makeatletter
\AtBeginDocument{\let\ps@plain\ps@nlrplain}
\makeatother

% The front matter in tex/before-body.tex is the title page. Pandoc's
% \maketitle would print a second, generic one ahead of it.
\renewcommand{\maketitle}{}

% --- 9. THE COMPONENT LIBRARY (spec B7) ------------------------------------
% Nothing here is a box any more.
%
% The old version of this section built every component the same way: a tinted
% panel, full measure, sharp corners, a coloured bar down the left edge, and
% for three of them an icon drawn in TikZ. That is a documentation-site
% admonition, and stacking three of them on one page is what made the book look
% machine-made no matter what typeface it was set in. The icons were the
% warning triangle, the green check and the circled X, which are the three most
% recognisable interface glyphs there are.
%
% What replaces them is the way print has always set an aside: a smaller size,
% an indent, and one italic word run into the first line. It costs no colour,
% survives a black-and-white printer, and cannot be mistaken for a screenshot.
\usepackage{fancyvrb}
\usepackage{pgffor}

% 9.1 Key terms ..............................................................
% There is no side column any more, so there is no margin note. A term is set
% bold where it first appears and links to its glossary entry. The definition
% lives in the glossary and nowhere else, which is one place rather than two
% that disagree.
\newcommand{\nlrnote}[1]{}
\newcommand{\nlrdef}[2]{\textbf{#1}}
\newcommand{\nlrnoteinline}[1]{}
\newcommand{\nlrdefinline}[2]{\textbf{#1}}
\newcommand{\nlrnomargin}{}

% The shared shape: indent, hold the indent for as long as the aside runs, and
% run an italic label into the first line.
% The indent is kept in a length rather than written into \leftskip directly,
% because anything inside an aside that sizes itself off \linewidth needs to
% subtract it. \linewidth does not know about \leftskip, so the error/fix pair
% laid two minipages across the full measure starting 1.4 em in, and ran 14
% lines past the right edge of the text block.
\newlength{\nlrasideindent}
\setlength{\nlrasideindent}{0pt}

\newenvironment{nlraside}[1]
  {\par\addvspace{11pt}\begingroup
   % 1.2em, resolved at body size before the callout size takes effect, so an
   % aside starts on the same vertical as an indented paragraph rather than 2pt
   % off it.
   \setlength{\nlrasideindent}{1.2em}%
   \setlength{\leftskip}{\nlrasideindent}\setlength{\parindent}{0pt}\nlrcalloutsize
   \noindent{\itshape #1.}\hspace{0.45em}\ignorespaces}
  {\par\endgroup\addvspace{11pt}}

% 9.2 Code block .............................................................
% Indented, in one ink, on the white page. That is how a printed listing has
% always been set. Never breaks across a page: a split code block is unreadable,
% so tools/check_layout.py flags any block long enough to risk it.
% Code sits flush with the text margin while prose is indented 1.2 em, so the
% two are told apart by the code hanging outward rather than by a panel. It has
% to be flush: the book guarantees 80 columns, and 80 columns of Inconsolata at
% 9 pt is 5.00 in against a 5.05 in measure, so any indent at all pushes the
% longest lines past the right edge of the text block.
%
% Inside an aside the whole block is indented, and \nlrasideindent carries that
% in so a listing there lines up with the prose around it.
\fvset{fontsize=\nlrcodesize, xleftmargin=\nlrasideindent, xrightmargin=0pt,
       samepage=false}

\newenvironment{nlrcodebox}[1]
  {\par\addvspace{9pt}%
   \if\relax\detokenize{#1}\relax\else
     \noindent\hspace{\nlrasideindent}{\nlrlabelmed{#1}}\par\vspace{1pt}%
   \fi}
  {\par\addvspace{9pt}}

% 9.3 Output block ...........................................................
% Beginners confuse input with output constantly, so the label is never
% optional. It used to be a grey panel under a dashed rule, which is a CSS
% border and a UI card; the distinction it was drawing now comes from the word
% "Output" and from the output being set in the secondary grey.
%
% This environment opened with \vspace{-10pt}, a kern tuned for the case where
% it follows a code block. Where it followed a paragraph instead, it dragged
% the block up over live text and the fill painted out a line: chapter 4 printed
% "The answer lands directly" and simply lost "underneath the box you typed in."
% \addvspace cannot do that.
\newenvironment{nlroutputbox}
  {\par\addvspace{4pt}%
   \noindent\hspace{\nlrasideindent}{\nlrlabelmed{Output}}\par\vspace{1pt}%
   \color{nlrsecond}\nlroutsize}
  {\par\addvspace{9pt}}

% 9.4 Exercise block .........................................................
% The one component that still announces itself, because a reader flipping for
% somewhere to start has to be able to find one. It does it with size and space
% rather than with a rule or a panel.
\newcommand{\nlrfield}[1]{%
  \par\addvspace{4pt}\noindent{\itshape #1.}\hspace{0.45em}\ignorespaces}

\newenvironment{nlrexercise}[2]
  {\par\addvspace{17pt}\begingroup\setlength{\parindent}{0pt}%
   \noindent{\bfseries\fontsize{13}{17}\selectfont #1}%
   \hfill{\itshape\color{nlrsecond}\fontsize{9.5}{13}\selectfont #2}%
   \par\addvspace{5pt}}
  {\par\endgroup\addvspace{17pt}}

% 9.5 "If it breaks" .........................................................
\newenvironment{nlrbreaks}
  {\begin{nlraside}{If it breaks}}
  {\end{nlraside}}

% One error/fix pair. The error message is set in mono so the reader can
% pattern-match it against what is on their own screen.
\usepackage{tabularx}
\newlength{\nlrerrwidth}
\newcommand{\nlrerrfix}[2]{%
  \par\addvspace{3pt}%
  \setlength{\nlrerrwidth}{\dimexpr\linewidth-\nlrasideindent\relax}%
  \noindent\begin{minipage}[t]{0.40\nlrerrwidth}\RaggedRight
    {\nlrmono\fontsize{8.5}{12}\selectfont #1}\end{minipage}%
  \hspace{0.02\nlrerrwidth}%
  \begin{minipage}[t]{0.56\nlrerrwidth}\RaggedRight\nlrcalloutsize #2\end{minipage}%
  \par\addvspace{3pt}}

% 9.5b "Trap" ...............................................................
% Distinct from "If it breaks" on purpose. An error stops you; a trap does not.
% This is for the case that runs cleanly and gives the wrong answer, which is
% the failure a beginner cannot detect. The two are told apart by the word,
% which is the only thing that ever told them apart in greyscale anyway.
\newenvironment{nlrtrap}
  {\begin{nlraside}{Trap}}
  {\end{nlraside}}

% 9.5c "Check" ..............................................................
\newenvironment{nlrcheck}
  {\begin{nlraside}{Check}}
  {\end{nlraside}}

% 9.5d Version pin ..........................................................
% One line. Web tools get redesigned and printed click-paths rot; this states
% when the page was last checked against the live service.
\newcommand{\nlrasof}[1]{%
  \par\addvspace{7pt}%
  \noindent{\nlrlabelmed{As of}\hspace{0.45em}%
  {\color{nlrsecond}\fontsize{9}{12}\selectfont #1}}%
  \par\addvspace{7pt}}

% 9.6 "Why this matters" ....................................................
\newenvironment{nlrwhy}
  {\begin{nlraside}{Why this matters}}
  {\end{nlraside}}

\newenvironment{nlrnotebox}[1]
  {\begin{nlraside}{#1}}
  {\end{nlraside}}

% 9.7 Checkpoint .............................................................
% Was a full-measure hairline with nothing closing it underneath, which is an
% <hr> followed by a <div>. The rule is gone; space and the italic line do it.
\newenvironment{nlrcheckpoint}
  {\par\addvspace{16pt}\begingroup\setlength{\parindent}{0pt}\nlrcalloutsize
   \noindent{\itshape You should now be able to:}\par\vspace{3pt}%
   \begin{itemize}[leftmargin=1.4em, rightmargin=0pt, itemsep=2pt,
                   topsep=0pt, parsep=0pt]}
  {\end{itemize}\endgroup\par\addvspace{14pt}}

% 9.8 Answer space ...........................................................
% These two stay ruled. They are not decoration: they are the place the reader
% writes, and a line to write on is the whole point of a workbook.
% 0.30 in rule spacing is a comfortable handwriting height for an adult hand.
\newcommand{\nlranswerlines}[1]{%
  \par\addvspace{6pt}%
  \foreach \i in {1,...,#1}{%
    \vspace{0.30in}\noindent{\color{nlrrule}\rule{\linewidth}{0.5pt}}\par}%
  \par\addvspace{8pt}}

\newcommand{\nlranswerbox}[1]{%
  \par\addvspace{6pt}%
  \noindent{\color{nlrrule}\fboxrule=0.5pt\fboxsep=8pt
    \fbox{\parbox[t][#1][t]{\dimexpr\linewidth-2\fboxsep-2\fboxrule\relax}{\strut}}}%
  \par\addvspace{8pt}}

% --- 10. LISTS, TABLES, FIGURES (spec B8) ----------------------------------
\usepackage{enumitem}
\setlist{leftmargin=1.2em, topsep=4pt, itemsep=3pt, parsep=0pt}
% Glossary entries. The stock description environment hangs its term into the
% margin, which puts it outside the text block, and it is handled by the block
% package rather than enumitem so its keys cannot be reconfigured. The filter
% converts the glossary's definition list into these two instead.
\newcommand{\nlrglossterm}[1]{%
  \par\addvspace{7pt}\noindent\textbf{#1}\par\nobreak\vspace{1pt}}
% \parindent has to be zeroed here, and the block indent has to be deeper than
% it was. At 1.1em the block sat *inside* the 1.2em first-line indent, so every
% glossary entry had three left edges and the first line of each definition
% started outside the block it belonged to.
\newenvironment{nlrglossdef}
  {\par\begingroup\setlength{\leftskip}{1.5em}\setlength{\parindent}{0pt}}
  {\par\endgroup\addvspace{2pt}}
% align=left keeps the number at the margin. The default right-aligns it into
% a fixed label box, which hangs it a fraction outside the measure.
% Markers hang into the margin, which is where LaTeX puts them and where books
% have always put them. The alternative is indenting every list body, which
% costs measure and makes a bulleted page look like a nested outline. Only
% `label` is set here: enumerate and itemize are handled by the block package
% in current LaTeX, and it rejects enumitem's geometry keys.
% Markers print in the text colour. Colouring the marker but leaving the text
% black is the CSS ::marker trick, and at 11 pt the oxblood bullet rendered as
% a faint red speck rather than as a mark.
\setlist[enumerate,1]{label=\arabic*.}
\setlist[itemize,1]{label=\textbullet}

\usepackage{booktabs}
\usepackage{longtable}
\usepackage{array}
\setlength{\aboverulesep}{3pt}
\setlength{\belowrulesep}{3pt}
\renewcommand{\arraystretch}{1.25}
% \nlrtablesize was defined and never used, so table cells were set at the body's
% 11/16. In the error index that put a wrapped line 24px below its own first line
% and the next entry only 30px below, and seven of thirteen rows on a page wrap,
% so their continuations read as separate entries. At 9.5/13 the two are far
% enough apart to tell.
\usepackage{etoolbox}
\AtBeginEnvironment{longtable}{\nlrtablesize}
\AtBeginEnvironment{tabular}{\nlrtablesize}

\usepackage{graphicx}
\usepackage{float}
% Captions are the second thing the tagging engine takes over. latex-lab's
% float module redefines \@makecaption inside a begindocument hook, so every
% \captionsetup key here was discarded and captions printed centred, at body
% size, in the roman: the one piece of the page that contradicted the flush-left
% ragged-right block it sat in. Overriding the socket and the macro is the
% route that survives, and it has to happen after their hook, not before it.
\makeatletter
\NewSocketPlug{caption/label}{nlr}{{\itshape #1.}}
\AssignSocketPlug{caption/label}{nlr}
\AddToHook{begindocument/end}{%
  \long\def\@makecaption#1#2{%
    \UseHookWithArguments{cmd/@makecaption/before}{2}{#1}{#2}%
    \vskip\abovecaptionskip
    \UseTaggingSocket{caption/begin}{\@current@float@struct}%
    {\nlrcapsize\color{nlrsecond}\RaggedRight\setlength{\parindent}{0pt}%
     \UseTaggingSocket{caption/label/begin}%
     \UseSocket{caption/label}{#1}%
     \UseTaggingSocket{caption/label/end}%
     \hspace{0.45em}%
     % \par closes the para structure by itself. Emitting para/end as well
     % double-closes it, and tagpdf refuses the document: 'The number of
     % automatic begin (3845) and end (3850) text para hooks differ'.
     \UseTaggingSocket{para/begin}#2\par}%
    \UseTaggingSocket{caption/end}%
    \vskip\belowcaptionskip}%
}
\makeatother
\setlength{\abovecaptionskip}{7pt}
\setlength{\belowcaptionskip}{2pt}

% --- 11. LINKS AND BOOKMARKS (spec B9) -------------------------------------
% Long addresses have to be able to break, or a single URL in the dataset index
% runs past the measure and there is nothing TeX can do about it.
\usepackage{xurl}
\usepackage{bookmark}
% Re-issued at the end of the preamble. Pandoc emits its own \hypersetup from
% the `colorlinks` metadata before this file is read, and at that point nlrteal
% does not exist yet, so every link fell back to LaTeX's default pure blue.
% On a printed copyright page that blue is the loudest thing on the sheet.
% Internal links print in the text colour. Coloured, they made every glossary
% term a bold oxblood phrase mid-paragraph and turned the error index into a
% column of ten red numerals -- so the same colour meant "defined term" here,
% "cross-reference" there and "page number" ten times over, and the only marks
% on a printed index page were red. An external address is the one link a
% reader may want to see as one.
\AtBeginDocument{\hypersetup{
  colorlinks = true,
  linkcolor  = nlrink,
  citecolor  = nlrink,
  filecolor  = nlrink,
  urlcolor   = nlraccent,
}}
\hypersetup{
  bookmarksnumbered = true,
  bookmarksopen     = true,
  bookmarksopenlevel = 1,
  pdfdisplaydoctitle = true,
}
\urlstyle{same}

% --- 12. QUARTO / PANDOC OVERRIDES -----------------------------------------
% Quarto wraps code in Shaded; the component library supplies the frame, so
% Shaded is flattened to a no-op here.
\makeatletter
\AtBeginDocument{%
  \@ifundefined{Shaded}{\newenvironment{Shaded}{}{}}{\renewenvironment{Shaded}{}{}}%
  \renewcommand{\verbatim@font}{\nlrmono\nlrcodesize}%
}
\makeatother

% Inline code: mono, marginally smaller. The Lua filter inserts \nlrbreak
% after separators so a long accession or path can wrap instead of hanging past
% the measure. A zero-width penalty, not a hyphen: a broken accession must not
% acquire punctuation that is not in it.
\newcommand{\nlrbreak}{\discretionary{}{}{}}
\let\nlroldtexttt\texttt
% 10 pt was too small: Inconsolata's x-height is shorter than Charis's, so the
% compound difference made every inline run visibly dip out of the line.
% No second argument: forcing 16 pt leading from inside \texttt pushed every
% table row that contained inline code to the body's leading, which inflated
% the error index until a wrapped line and a new entry were 3 pt apart.
\makeatletter
\renewcommand{\texttt}[1]{{\nlrmono\fontsize{10.4}{\f@baselineskip}\selectfont #1}}
\makeatother
